\documentclass[11pt,parskip=half]{scrartcl} % Lädt die Einstellungen der Artikel-Dokumentklasse des KOMA-Scripts
\usepackage[utf8]{inputenc} % Ermöglicht die Eingabe von deutschen Sonderzeichen im Quelltext
\usepackage[T1]{fontenc} % Europäische Kodierung der Ausgabeschrift
\usepackage{lmodern} % Lädt die Schriftart Latin Modern, die universell skalierbar ist
\usepackage[ngerman]{babel} % Silbentrennung nach neuer deutscher Rechtschreibung
\usepackage{amsmath} % Mathematik-Paket der American Mathematical Society
\usepackage{graphicx} % Erlaubt das Einfügen von Bildern per \includegraphics{}

\begin{document}
% Titel-Sektion
\author{Marina Mustermann, Bob Beispiel}
\title{Titel der Projektarbeit}
\subtitle{Untertitel der Projektarbeit}
\date{\today}
\maketitle

% Inhaltsverzeichnis
\tableofcontents

\section{Einleitung}
In der Einleitung wird in das Thema eingeführt und die Fragestellung beschrieben.

\section{Hauptteil}
Im Hauptteil wird die Fragestellung in einzelnen Abschnitten abgearbeitet.
\subsection{Grundlagen}
Mathematische Arbeiten beginnen meist mit einem Abschnitt über die Grundlagen des Themas. Hier können Begriffe, die für den Rest der Arbeit benötigt werden, definiert und erklärt werden.
\subsection{Extrem komplizierte mathematische Überlegungen}
Die Ableitung $f'$ einer Funktion $f$ ist definiert als Grenzwert des Differenzenquotienten für $h \rightarrow 0$, das heißt:
\begin{equation}
	f'(x)=\lim_{h \rightarrow 0} \frac{f(x+h)-f(x)}{h}.
\end{equation}

\section{Fazit und Ausblick}
In einem Fazit wird das Ergebnis der Projekt-Arbeit zusammengefasst und beurteilt, inwiefern die gegebene Fragestellung beantwortet werden konnte. Wurde die Fragestellung nur teilweise beantwortet, so muss eine Diskussion der Ursachen stattfinden. Außerdem kann man einen Ausblick auf weitere Fragestellungen geben, die sich aus der Arbeit ergeben haben.

\section{Literaturverzeichnis}
Im Literaturverzeichnis werden die Quellen (Bücher, Internetseiten) mit Autor und Erscheinungsdatum alphabetisch sortiert aufgeführt. In der Arbeit wird dann auf die Quellen aus dem Literaturverzeichnis verwiesen.

\end{document} 
